\begin{problem}{}{}{}{}{}
\vspace{-4.0em}
\begin{center}
    {\LARGE\textbf{<Bài ...>. <Tên bài>} (<điểm nếu có>)}\\[0.7em]
    \textbf{Thời gian:} ... giây \\ % xóa dòng này nếu giới hạn thời gian không được đề cập
    \textbf{Bộ nhớ:} ... megabytes \\ % xóa dòng này nếu giới hạn bộ nhớ không được đề cập
    \textbf{Dữ liệu:} \texttt{<standard input hoặc tên file input>} \\
    \textbf{Kết quả:} \texttt{<standard output hoặc tên file output>}
\end{center}

<Phát biểu đề bài>

\medskip

\InputFile

Nhập từ tệp văn bản \texttt{<tên file input>}: % Nếu đề bài yêu cầu nhập từ file
% Nhập từ bàn phím: % Nếu đề bài yêu cầu nhập từ stdin

\begin{itemize}
    \item <Định dạng input kèm ràng buộc> % Ví dụ: Dòng đầu tiên gồm một số nguyên dương $n$ $(1 \le n \le 10^5)$.
\end{itemize}

\OutputFile

In ra tệp văn bản \texttt{<tên file output>}: % Nếu đề bài yêu cầu nhập từ file
% In ra màn hình: % Nếu đề bài yêu cầu in ra stdout

\begin{itemize}
    \item <Định dạng output>
\end{itemize}

\Subtask{}
\begin{itemize}
    \item Subtask ...: $...\%$ số điểm có ....
\end{itemize}

\medskip

\Examples
\begin{center}
\begin{example}
\exmp{
<sample input 1>
}{
<sample output 1>
}%
\exmp{
<sample input 2>
}{
<sample output 2>
}%
\end{example}
\end{center}

% Sau mỗi test ví dụ phải có ký tự phần trăm (%) để không khiến cho khung test ví dụ bị thừa.

\Note

<Giải thích ví dụ>

% Lưu ý không ghi lại từ "Giải thích"

\end{problem}